% Personal Style LLM — systems / case-study paper (ACL format).
%
% Compiles two ways (see paper/README.md):
%   * With the ACL Rolling Review template: drop acl.sty + acl_natbib.bst into
%     this directory (or use the Overleaf template) and it builds in ACL style.
%   * Without it: falls back to a plain two-column article so the draft always
%     compiles. The fallback is NOT camera-ready format.
%
% Numbers not yet measured are wrapped in \pending{} so an unfilled draft is
% visibly unfinished rather than quietly wrong. Every \pending must be replaced
% by a value traceable to a row under artifacts/eval/ before submission.

\IfFileExists{acl.sty}{%
  \documentclass[11pt]{article}
  \usepackage[review]{acl}
}{%
  \documentclass[11pt,twocolumn]{article}
  \usepackage[margin=2.5cm]{geometry}
  \usepackage{natbib}
  \bibliographystyle{plainnat}
  \PassOptionsToPackage{hyphens}{url}
  \usepackage{hyperref}
}

\usepackage{times}
\usepackage{latexsym}
\usepackage{booktabs}
\usepackage{graphicx}
\usepackage{amsmath}
\usepackage{xcolor}
\usepackage[T1]{fontenc}
\usepackage[utf8]{inputenc}
\usepackage{microtype}

% Visibly-unfilled result placeholder.
\newcommand{\pending}[1][??]{\textcolor{red}{\textbf{#1}}}

\title{Writing in One Person's Voice: A Case Study in Local,\\
       Retrieval-Conditioned Personal Style Generation}

\author{Shiloh Siu \\
  \texttt{y.shing@tranxfer.com} \\}

\begin{document}
\maketitle

\begin{abstract}
We report a case study in generating text in a single author's personal writing
voice, under a hard constraint that no text ever leaves the author's machine.
The system profiles a personal corpus along lexical, syntactic and pragmatic
dimensions, retrieves diverse in-genre exemplars from a local dense index, and
conditions a local instruct model on both. We evaluate it two ways: a
stylometric similarity metric over \pending{648} generations spanning three
open-weight models, six document genres and a four-rung ablation ladder; and a
blind A/B protocol in which the author, the only person who can judge the
target style, rates styled against unstyled drafts. Measurement also corrected
the project's own premise: the corpus the system was designed around is
\emph{smaller} than assumed, because the document projected to dominate it
turned out to contribute no extractable prose. We report the ablation result
(\pending), the human preference rate (\pending\%, $n=\pending{20}$), and what a
proxy stylometric metric can and cannot establish at $n{=}1$ author.
\end{abstract}

\section{Introduction}
\label{sec:intro}

Large instruct models write competent generic prose. They do not, by default,
write like any particular person. For an individual who wants drafting help with
the documents they actually produce --- personal statements, cover letters,
coursework, reflective journals --- the gap that matters is not fluency but
voice.

Closing that gap with fine-tuning is the obvious move and, for a single author,
the wrong first one: the corpus is small enough that fine-tuning risks
memorisation rather than abstraction, and the compute and data-handling
requirements sit badly with the privacy constraint that motivates the project.
This paper instead asks how far \emph{prompt-time conditioning} gets, and
measures it.

The setting is deliberately narrow. One author, one corpus of 17 personal
documents, everything on one machine. That narrowness is the point of a case
study: it lets us report what actually happened, including the parts that
contradicted the design assumptions, rather than an averaged result over a
population that no individual user resembles.

\paragraph{Contributions.}
\begin{enumerate}\itemsep0pt
  \item A local, offline pipeline for personal-style generation: skew-aware
        style profiling, dense retrieval with MMR diversity, and prompt
        orchestration over an on-device open-weight model
        (\S\ref{sec:system}).
  \item A stylometric similarity metric appropriate to the task, and an argument
        for why the conventional choices (BLEU, embedding cosine) measure the
        wrong thing here (\S\ref{sec:metric}).
  \item An ablation over the two conditioning components, across three models and
        six genres, plus a blind human A/B (\S\ref{sec:results}).
  \item A negative measurement result about our own corpus that overturned the
        design's central assumption (\S\ref{sec:corpus}).
\end{enumerate}

\section{Related Work}
\label{sec:related}

\paragraph{Stylometry and authorship.}
Quantifying authorial style through function words, sentence-length
distributions and part-of-speech patterns is long-established in authorship
attribution. We borrow its feature vocabulary, but invert the task: attribution
asks whether a text matches an author, while we ask a generator to move toward
that match.

\paragraph{Controllable and style-conditioned generation.}
Prior work conditions generation on style through fine-tuning, control codes, or
in-context exemplars. The single-author, small-corpus regime pushes toward
in-context methods for the reasons given in \S\ref{sec:intro}.

\paragraph{Retrieval-augmented generation.}
RAG is normally used to supply \emph{facts}. Here retrieval supplies \emph{form}:
the retrieved chunks are not evidence to be reproduced but demonstrations of how
the author writes. This changes what good retrieval means --- diversity of
rhetorical register matters more than topical precision, which is why selection
uses MMR (\S\ref{sec:retrieval}).

\paragraph{Why not BLEU or embedding cosine.}
$n$-gram overlap with the author's existing documents rewards a generator for
copying them, which is the failure mode we most want to avoid. Sentence-embedding
cosine is dominated by topic: a cover letter and a physics investigation by the
same author are far apart in embedding space despite sharing a voice. Both
metrics would report progress on something other than style
(\S\ref{sec:metric}).

\section{System}
\label{sec:system}

\begin{figure}[t]
\centering
\small
\texttt{PDFs} $\rightarrow$ \texttt{ingest} $\rightarrow$ \texttt{style profile} \\[2pt]
$\downarrow$ \\[2pt]
\texttt{chunk + embed} $\rightarrow$ \texttt{dense index} \\[2pt]
$\downarrow$ \\[2pt]
\texttt{retrieve (MMR)} $\rightarrow$ \texttt{prompt assembly} $\rightarrow$ \texttt{local LLM} \\[2pt]
$\downarrow$ \\[2pt]
\texttt{stylometric score} \;/\; \texttt{blind A/B}
\caption{Pipeline. Every stage runs on the author's machine; model weights are
fetched once and inference is on-device.}
\label{fig:pipeline}
\end{figure}

\subsection{Ingestion and skew-aware profiling}
\label{sec:profiling}

Documents are extracted from PDF, stripped of running headers, footers,
bibliographies and non-prose lines, tagged with a genre label, and screened by a
soft PII pass. Style features are then computed along four dimensions: lexical
(type-token ratio, MTLD), syntactic (mean sentence length, POS distribution),
pragmatic (formality, reading grade, rhetorical markers) and semantic
(topic clusters, embedding centroid).

The profile is computed \emph{per genre and then averaged across genres} rather
than pooled over the corpus. This was a defensive choice: the corpus was
projected to be dominated by one very long document, and pooling would have let
that document's register stand in for the author's. \S\ref{sec:corpus} reports
what measurement did to that assumption.

\subsection{Retrieval}
\label{sec:retrieval}

Chunks are paragraph-packed to a token target, never split mid-sentence, and
embedded with a sentence-transformer model. Retrieval is brute-force cosine over
the flat index --- justified at this corpus size --- with an optional genre
filter, near-duplicate collapsing, and MMR selection so the top-$k$ exemplars
demonstrate a range of the author's constructions rather than five near-copies of
one paragraph. No BM25 and no cross-encoder reranker: neither is defensible at a
few dozen chunks.

\subsection{Prompt orchestration}
\label{sec:prompting}

The assembled prompt carries a system instruction, a natural-language
\emph{style brief} rendered from the profile, the retrieved exemplars labelled as
models of tone rather than content, and the request. The two conditioning
components are independently switchable, which is exactly what the ablation in
\S\ref{sec:setup} varies.

Generation runs locally through Hugging Face \texttt{transformers}: weights are
downloaded once from the Hub, after which runs execute with the hub client in
offline mode. Chain-of-thought is suppressed on hybrid-reasoning models, and any
residual reasoning span is stripped before scoring --- left in, it would corrupt
the stylometric vector with text the author never intended as prose.

\subsection{The stylometric similarity metric}
\label{sec:metric}

A draft is scored by extracting the same feature vector used for profiling and
measuring its distance to the author's target vector for that genre. Features
differ in natural magnitude by orders of magnitude, so each is standardised by
its per-feature standard deviation across the corpus before the Euclidean
distance is taken; the distance $d$ is mapped to a similarity by $1/(1+d)$.

We stress that this is a \emph{proxy}. It measures whether a draft occupies the
author's region of stylometric feature space. A text can do that while reading
as obviously machine-written, and the metric would not notice. It is used here
because it scales to hundreds of generations and is sensitive to the right
\emph{kind} of difference; the human A/B is what tests whether the result is
actually convincing.

\section{Corpus}
\label{sec:corpus}

The corpus is 17 PDF documents written by the author across six genres:
personal statements, cover letters, academic investigations, coursework
assignments, reflective journals, and one extended essay.

\paragraph{A premise overturned by measurement.}
The system was designed around a projected corpus of roughly 93{,}900 words in
which a single 395-page extended essay accounted for about 69\% of the total ---
a skew severe enough to justify the per-genre weighting in
\S\ref{sec:profiling}. Ingestion showed this to be wrong. The extended essay's
PDF text layer contains only its numeric data appendix; the essay prose itself is
not extractable from the file. After the minimum-token screen the document
contributes nothing to either the profile or the index.

The real corpus is therefore \pending{16} usable documents yielding \pending{56}
chunks --- roughly a fifth of the projected size. Two consequences follow, and we
report both rather than quietly re-scoping. First, the skew-aware design is
largely moot: with the dominant document absent, per-genre averaging and naive
pooling are much closer than anticipated, so the defensive choice cost little but
bought little. Second, the \texttt{extended\_essay} genre has no exemplars at
all, so its exemplar-bearing ablation conditions collapse onto the
profile-only ones. Results for that genre are reported separately and should not
be read as a retrieval result.

The general lesson is unremarkable but easy to skip: a corpus assumption stated in
a design document is a hypothesis, and ingestion is the experiment that tests it.

\section{Experimental Setup}
\label{sec:setup}

\paragraph{Models.} Three open-weight instruct models in the 7--9B class, run
locally (Table~\ref{tab:models}). Commit revisions are pinned so the campaign is
reproducible against a moving model hub.

\begin{table}[t]
\centering\small
\begin{tabular}{@{}lll@{}}
\toprule
Key & Repository & Revision \\
\midrule
\texttt{qwen2.5-7b}  & Qwen/Qwen2.5-7B-Instruct        & \pending \\
\texttt{llama3.1-8b} & meta-llama/Llama-3.1-8B-Instruct & \pending \\
\texttt{mistral-7b}  & mistralai/Mistral-7B-Instruct-v0.3 & \pending \\
\bottomrule
\end{tabular}
\caption{Models. Hardware: \pending{GPU}, \pending{VRAM}.}
\label{tab:models}
\end{table}

\paragraph{Prompt bank.} 18 fixed generation requests, three per genre, written
as realistic requests in that genre and versioned in the repository so results
remain interpretable.

\paragraph{Ablation ladder.} Four conditions crossing the two conditioning
components: \texttt{plain} (neither), \texttt{profile\_only} (style brief),
\texttt{exemplars\_only} (retrieved exemplars), \texttt{full} (both).

\paragraph{Repetitions and seeding.} Decoding samples at temperature $0.7$, so
each cell is run three times. The sampler is seeded per repetition, making each
(cell, repetition) individually reproducible while still exposing run-to-run
variance --- $3 \times 18 \times 4 \times 3 = 648$ generations.

\paragraph{A/B protocol.} \pending{20} prompts sampled balanced across genres.
For each, a styled and a plain draft are generated and presented under randomised
blind labels; the author records which reads more like their own writing. The
author is the only available rater --- the target style is theirs --- and this is
a real limitation, treated as such in \S\ref{sec:limitations}.

\section{Results}
\label{sec:results}

\begin{table}[t]
\centering\small
\begin{tabular}{@{}lrrr@{}}
\toprule
Condition & $n$ & Mean sim. & SD \\
\midrule
\texttt{plain}          & \pending & \pending & \pending \\
\texttt{profile\_only}   & \pending & \pending & \pending \\
\texttt{exemplars\_only} & \pending & \pending & \pending \\
\texttt{full}           & \pending & \pending & \pending \\
\bottomrule
\end{tabular}
\caption{Ablation ladder, pooled over models and genres. Source:
\texttt{artifacts/eval/table\_by\_condition.csv}.}
\label{tab:ablation}
\end{table}

\begin{figure}[t]
\centering
\includegraphics[width=\columnwidth]{figures/fig_ablation.pdf}
\caption{Mean stylometric similarity per condition, grouped by model. Error bars
are one standard deviation.}
\label{fig:ablation}
\end{figure}

\begin{figure}[t]
\centering
\includegraphics[width=\columnwidth]{figures/fig_by_doc_type.pdf}
\caption{The same ladder broken out by genre.}
\label{fig:bydoctype}
\end{figure}

\paragraph{Ablation.} \pending{Describe which rungs separate.} Conditions are
compared pairwise against \texttt{full} with a Wilcoxon signed-rank test, paired
by (model, prompt, repetition) so the only difference within a pair is what the
prompt carried; the test is used in place of a $t$-test because the differences
are small-sample and not plausibly normal. Source:
\texttt{artifacts/eval/ablation\_tests.csv}.

\paragraph{Models and genres.} \pending{Per-model and per-genre findings.}

\paragraph{Human A/B.} The styled arm was preferred in \pending{}/\pending{20}
trials (\pending\%, 95\% Wilson CI [\pending, \pending], binomial $p =
\pending$ against a 0.5 null). \pending{Interpretation.}

\paragraph{Latency.} \pending{Mean latency and tokens/s per model.}

\section{Discussion}
\label{sec:discussion}

\pending{To be written from the measured results.} The questions the design
poses, and which the results should answer, are: whether the style brief and the
retrieved exemplars contribute independently or redundantly; whether either
component helps more in genres with strong conventional form (cover letters) than
in genres where the author's voice carries more of the work (reflections); and
whether the automated metric and the human preference agree --- a disagreement
would be the more informative outcome, since it would bound what the proxy can be
trusted to say.

\section{Limitations}
\label{sec:limitations}

\paragraph{One author.} Every result is about one person's writing. Nothing here
establishes that the approach generalises, and the corpus-size finding in
\S\ref{sec:corpus} suggests personal corpora vary in ways that matter.

\paragraph{One rater, who is also the author of the system.} The A/B judgments
come from the same person who built the pipeline and whose style is the target.
Only that person can judge the target, but they are not blind to the study's
purpose, and at $n \approx 20$ the confidence interval is wide. The preference
rate should be read as evidence about direction, not magnitude.

\paragraph{A proxy metric.} The stylometric similarity measures position in
feature space, not whether prose is convincing. It cannot detect a draft that is
stylometrically on-target and obviously synthetic.

\paragraph{Small corpus.} \pending{16} documents and \pending{56} chunks. The
per-genre profile for thinly represented genres rests on very few documents, and
one genre has no usable prose at all.

\paragraph{Stochastic decoding.} Seeding makes each run reproducible but does not
remove sampling variance; three repetitions bound it only loosely.

\paragraph{No fine-tuning baseline.} The paper measures how far prompt-time
conditioning gets, without establishing what fine-tuning would add. That
comparison is the natural next study, not a result reported here.

\section{Ethics and Privacy}
\label{sec:ethics}

The corpus is one person's personal writing, including university applications.
It is used with that person's consent --- they are the author of both the corpus
and this paper --- and it is never transmitted: ingestion, indexing, retrieval,
generation and scoring all run on one machine, model weights are fetched once
from the model hub and inference is on-device, and no hosted inference API is
used at any point. The source corpus and all derived artefacts are excluded from
version control. A soft PII pass flags obvious emails and phone numbers during
ingestion; it is a convenience, not a safety guarantee, and is not relied on as
one.

A system that writes convincingly in a named person's voice is dual-use. The
mitigation here is structural rather than technical: it runs locally, on a
corpus its own subject supplies, with no capability to acquire someone else's
writing.

\section{Conclusion}
\label{sec:conclusion}

\pending{To be written once the A/B bar is measured.} The decision this study
exists to inform is whether prompt-time conditioning is sufficient, or whether a
second version should invest in fine-tuning or hybrid retrieval. That decision
should follow from the measured preference rate and its interval, not from the
proxy metric alone.

\bibliography{references}

\appendix

\section{Prompt Bank}
\label{app:prompts}
The 18 requests are listed in \texttt{experiments/prompts.yaml}.

\section{Full Result Tables}
\label{app:tables}
Per-model, per-genre and per-condition tables are generated into
\texttt{artifacts/eval/} by \texttt{experiments/analyze\_results.py}.

\section{Example Generations}
\label{app:examples}
\pending{Selected drafts per condition.}

\end{document}
